\rok{Новембар 2024.}{H}

Први практични тест из Алгоритама и структура података

\zadatak{1}{Merge Sort}
Написати методу \textit{merge()} \textit{MergeSort} алгоритма за сортирање целобројног низа дужине \textit{n}.

\textbf{Додатне напомене за решавање задатка:}
\begin{itemize}
    \item Улаз је несортиран низ целих бројева.
    \item Излаз је сортиран низ целих бројева.
    \item У шаблону се налази класа \texttt{RandomNumberGenerator} коју треба искористити за генерисање низа случајних бројева.
    \item У оквиру класе \texttt{MergeSort} написати имплементацију за методу:
    \begin{Verbatim}[fontsize=\footnotesize, numbers=left, xleftmargin=10mm]
private static void merge(long[] arr, long[] tmpArray,
                int lowerIndex, int middleIndex, int upperIndex)
    \end{Verbatim}
    \item Обезбедити код у \texttt{Main} методи за тестирање алгоритма.
\end{itemize}



\zadatak{2}{Промена редоследа \textit{n} елемената реда, коришћењем стека}
Имплементирати функцију која за задати ред целих бројева обрће редослед само првих \textit{n} елемената реда, користећи помоћни стек. У оквиру тих првих \textit{n} елемената, за сваки од њих који је паран важиће да се његова вредност повећава за 2. Након тога, преостали елементи реда остају у истом редоследу.

\textbf{Спецификација функције:}
\begin{itemize}
    \item Функција прима број \textit{n}, где је \textit{n} број елемената које треба обрнути.
    \item Функција треба користити стек (\textit{stack}) као помоћну структуру.
    \item Елементи после првих \textit{n} треба да остану у оригиналном редоследу.
\end{itemize}

\textbf{Додатни захтеви за решавање задатка:}
\begin{itemize}
    \item Алгоритам писати у оквиру класе \texttt{Queue}.
    \item Захтеван потпис методе:
    \begin{Verbatim}[fontsize=\footnotesize, numbers=left, xleftmargin=10mm]
public void ReverseTopNQueueElements(int n)
    \end{Verbatim}
\end{itemize}

\textbf{Примери:}
\begin{itemize}
    \item \textbf{Пример 1:} \\
          \textbf{Оригинални ред:} 1 2 3 4 5 6 7, $n = 4$ \\
          Међу прва 4 елемента, 2 и 4 јесу парни бројеви, те се њихова вредност повећава за 2. Зато у испису 4 постаје 6, а 2 постаје 4. \\
          \textbf{Ред након обртања:} 6 3 4 1 5 6 7
    \item \textbf{Пример 2:} \\
          \textbf{Оригинални ред:} 1 2 3 4 5 6 7, $n = 7$ \\
          \textbf{Ред након обртања:} 7 8 5 6 3 4 1
    \item \textbf{Пример 3:} \\
          \textbf{Оригинални ред:} 1 2 3 4 5 6 7, $n = 8$ \\
          Неважећа вредност за \textit{n}.
\end{itemize}



\zadatak{3}{Најдужи сегмент са растућим и опадајућим вредностима}
Чворови у повезаној листи могу формирати \textbf{сегмент са растућим и опадајућим вредностима} ако задовољавају следеће услове:
\begin{enumerate}
    \item Чворови унутар сегмента прво имају строго растуће вредности, а затим строго опадајуће вредности.
    \item Сегмент мора имати најмање 3 чвора: најмање 2 у растућем делу и 1 у опадајућем делу, или обрнуто.
\end{enumerate}

Задатак је да се пронађе и врати:
\begin{enumerate}
    \item Дужину најдужег сегмента са растућим и опадајућим вредностима.
    \item Почетну вредност сегмента.
    \item Крајњу вредност сегмента.
\end{enumerate}

Ако ниједан такав сегмент не постоји, треба вратити низ \texttt{[-1, -1, -1]}.

\textbf{Додатни захтеви:}
\begin{itemize}
    \item Решавати овај задатак помоћу структуре листе (\texttt{LinkedList}) из шаблона задатка.
    \item Имплементирати у оквиру методе:
    \begin{Verbatim}[fontsize=\footnotesize, numbers=left, xleftmargin=10mm]
public int[] findLongestSegment()
    \end{Verbatim}
\end{itemize}

\textbf{Примери:}
\begin{itemize}
    \item \textbf{Пример 1:} \\
          \textbf{Улаз:} \texttt{1 -> 2 -> 3 -> 2 -> 1} \\
          \textbf{Излаз:} \texttt{[5, 1, 1]} \\
          Објашњење: Најдужи сегмент са растућим и опадајућим вредностима има дужину 5 (\texttt{1 -> 2 -> 3 -> 2 -> 1}), почиње од 1 и завршава се на 1.
    \item \textbf{Пример 2:} \\
          \textbf{Улаз:} \texttt{5 -> 4 -> 3 -> 2 -> 1} \\
          \textbf{Излаз:} \texttt{[-1, -1, -1]} \\
          Објашњење: Нема растућег дела у листи, па не постоји валидан сегмент.
\end{itemize}

\sablon{AISP2024_Test1_GrupaH}{2024/Novembar/GrupaH/AISP2024_Test1_GrupaH.linq}
